\documentclass[11pt]{amsart}

\usepackage[margin=1.08in]{geometry}
\usepackage{amsmath,amssymb,amsthm}
\usepackage{enumitem}
\usepackage{graphicx}
\usepackage[hidelinks]{hyperref}

\newtheorem{theorem}{Theorem}
\newtheorem{lemma}{Lemma}
\theoremstyle{remark}
\newtheorem{remark}{Remark}

\newcommand{\R}{\mathbb R}
\newcommand{\Lin}{\mathcal L}
\newcommand{\co}{\operatorname{co}}
\newcommand{\dist}{\operatorname{dist}}
\newcommand{\CNA}{\operatorname{CNA}}

\title{A square counterexample to Crawford-number localization near convex generators}
\author{Agent lane 04}
\date{Run \texttt{fa\_banach\_001}, 2026-06-29}

\begin{document}
\maketitle

\section*{Status}

This packet gives a finite-dimensional counterexample to Question 3.2 in
Choi--Lee, \emph{On the Crawford number attaining operators} \cite{ChoiLee2022}.
The question asks whether, whenever $S\subset B_X$ and
$B_X=\overline{\co}(S)$, the Crawford number $c(T)$ can be computed by
restricting the infimum to states $(x,x^*)$ with $\dist(x,S)<\varepsilon$.
The answer is no, even for a two-dimensional real Banach space and a finite
uniformly strongly exposed generating set $S$.

\section*{Source question}

The source asks:
\begin{quote}
Let $X$ be a Banach space and $S\subset B_X$ be a subset such that
$B_X=\overline{\co}(S)$. Given $\varepsilon>0$ and $T\in\mathcal L(X)$, does
\[
 c(T)=\inf\{|x^*(Tx)|:(x,x^*)\in\Pi(X),\ \operatorname{dist}(x,S)<\varepsilon\}
\]
hold?
\end{quote}

\begin{center}
\includegraphics[width=0.94\textwidth]{figures/open_problem_crop.png}
\end{center}

\section*{Definitions}

For a Banach space $X$, the state space is
\[
        \Pi(X)=\{(x,x^*)\in S_X\times S_{X^*}:x^*(x)=1\}.
\]
For $T\in\Lin(X)$, the Crawford number is
\[
        c(T)=\inf\{|x^*(Tx)|:(x,x^*)\in\Pi(X)\}.
\]

\section*{Idea of the proof}

In the square $X=\ell_\infty^2$, the four vertices generate the unit ball and
are uniformly strongly exposed.  However, the middle point of an edge can
carry a state that sees the coordinate-swap operator as zero.  Points close to
the vertices have both coordinates close to signs $\pm1$, and every norming
state is supported on a coordinate where the point has modulus one.  After the
coordinate swap, that state sees the other coordinate, whose modulus is still
close to one.  Thus the global Crawford infimum is zero, but the infimum
restricted near the generating set is bounded below by a positive number.

\section*{Counterexample}

\begin{theorem}
Question 3.2 of \cite{ChoiLee2022} has a negative answer.  More precisely,
there are a real two-dimensional Banach space $X$, a finite uniformly strongly
exposed set $S\subset B_X$ with $B_X=\co(S)$, an operator $T\in\Lin(X)$, and
$0<\varepsilon<1$ such that
\[
        c(T)=0
\]
but
\[
 \inf\{|x^*(Tx)|:(x,x^*)\in\Pi(X),\ \dist(x,S)<\varepsilon\}
        \ge 1-\varepsilon>0.
\]
\end{theorem}

\begin{proof}
Let $X=(\R^2,\|\cdot\|_\infty)$ and let
\[
        S=\{(\sigma,\tau):\sigma,\tau\in\{-1,1\}\}
\]
be the four vertices of the unit square.  Then $B_X=\co(S)$.
The set $S$ is uniformly strongly exposed: the vertex
$v_{\sigma,\tau}=(\sigma,\tau)$ is exposed by
\[
        f_{\sigma,\tau}(a,b)=\frac{\sigma a+\tau b}{2}.
\]
If $f_{\sigma,\tau}(a,b)>1-\eta$ for $(a,b)\in B_X$, then
$\sigma a>1-2\eta$ and $\tau b>1-2\eta$, so
$\|(a,b)-v_{\sigma,\tau}\|_\infty<2\eta$.  Since $S$ is finite, this is a
uniform strong exposure.

Define the coordinate-swap operator
\[
        T(a,b)=(b,a).
\]
Let $x=(1,0)$ and $x^*=e_1^*$, the first coordinate functional.  Then
$(x,x^*)\in\Pi(X)$ and
\[
        x^*(Tx)=e_1^*(0,1)=0.
\]
Hence $c(T)=0$.

Fix $0<\varepsilon<1$.  Suppose $(x,x^*)\in\Pi(X)$ and
$\dist(x,S)<\varepsilon$.  Write $x=(a,b)$.  There are
$\sigma,\tau\in\{-1,1\}$ such that
\[
        |a-\sigma|<\varepsilon,\qquad |b-\tau|<\varepsilon.
\]
In particular $|a|>1-\varepsilon$ and $|b|>1-\varepsilon$.

Write $x^*=(\alpha,\beta)\in(\ell_\infty^2)^*=\ell_1^2$.  Since
$(x,x^*)\in\Pi(X)$,
\[
        |\alpha|+|\beta|=1,\qquad \alpha a+\beta b=1.
\]
The equality case in
\[
        \alpha a+\beta b
        \le |\alpha|\,|a|+|\beta|\,|b|
        \le |\alpha|+|\beta|=1
\]
implies that every coordinate on which $x^*$ is supported is a norming
coordinate of $x$, with matching sign.  Thus, if $\alpha\ne0$ then
$|a|=1$ and $\alpha a=|\alpha|$, while if $\beta\ne0$ then $|b|=1$ and
$\beta b=|\beta|$.

Now
\[
        x^*(Tx)=\alpha b+\beta a.
\]
If only the first coordinate is used by $x^*$, then
$|x^*(Tx)|=|b|>1-\varepsilon$.  If only the second is used, then
$|x^*(Tx)|=|a|>1-\varepsilon$.  If both are used, then $|a|=|b|=1$ and
the signs of $a$ and $b$ are $\sigma$ and $\tau$, respectively; in that case
\[
        x^*(Tx)=\sigma\tau(|\alpha|+|\beta|)
\]
and so $|x^*(Tx)|=1$.

Therefore every state with $\dist(x,S)<\varepsilon$ satisfies
$|x^*(Tx)|\ge 1-\varepsilon$.  The restricted infimum is positive, whereas
$c(T)=0$, proving the claimed failure.
\end{proof}

\section*{Limitations}

This counterexample answers Question 3.2 as stated, and even with $S$ finite
and uniformly strongly exposed.  It does not settle the separate questions in
the source about density of $\CNA(X)$ for property $\alpha$ or property
$\beta$ spaces.

\section*{Novelty and verification notes}

\begin{itemize}[leftmargin=*]
\item Local run indexes were searched for \texttt{2201.10031}, the source
title, \texttt{Crawford number}, \texttt{property alpha}, and
\texttt{CNA(X)}.  No duplicate packet or exact attempt for this question was
found.
\item Bounded web search on 2026-06-29 for the source title and phrases such
as ``Crawford number property alpha'' and ``Crawford number closed convex
hull'' surfaced the source arXiv page but no later answer to this localization
question.
\item The example is finite dimensional and uses only the elementary
description of states in $\ell_\infty^2$.
\end{itemize}

\section*{Human review recommendation}

Review as a likely-valid counterexample to Question 3.2.  The key point to
check is the elementary equality-case argument for states of $\ell_\infty^2$.
If accepted, the counterexample also shows that adding ``$S$ uniformly strongly
exposed'' is not sufficient to rescue the localization formula.

\begin{thebibliography}{9}

\bibitem{ChoiLee2022}
G. Choi and H. J. Lee,
\emph{On the Crawford number attaining operators},
arXiv:2201.10031, 2022.
\url{https://arxiv.org/abs/2201.10031}

\end{thebibliography}

\end{document}
